\section{Methods}
\label{sec:methods}

Our analyses comprises of the following steps:

\paragraph{Data Integration}

Substance records were compiled from a comprehensive set of 18 European regulatory and monitoring sources and one Flemish administrative list. These include ECHA obligation lists (Candidate List~\cite{ECHA_CandidateList}, Authorisation List~\cite{ECHA_AuthorisationList}, Restriction List~\cite{ECHA_RestrictionList}, POPs List~\cite{ECHA_POPsList}, EU Positive List~\cite{ECHA_EUPositiveList}, and the CLH Harmonised Classification and Labelling List~\cite{ECHA_CLHList}), ECHA activity lists (Restriction Process~\cite{ECHA_RestrictionProcess}, SVHC Identification~\cite{ECHA_SVHCIdentification}, Authorisation Process~\cite{ECHA_AuthorisationProcess}, Dossier Evaluation~\cite{ECHA_DossierEvaluation}, CLH Process~\cite{ECHA_CLHProcess}, Substance Evaluation~\cite{ECHA_SubstanceEvaluation}, POPs Process~\cite{ECHA_POPsProcess}, PBT Assessment~\cite{ECHA_PBTAssessment}, and ED Assessment~\cite{ECHA_EDAssessment}), as well as REACH registration data~\cite{ECHA_REACHRegistrations}. 
%Table~\ref{tab:datasources} provides an overview of all regulatory and monitoring sources included in this study.

These are the most autoritative datasets that lay the foundation of EU regulatory oversight of chemical Substances.

Additional sources comprise the EU Pesticides Database (Active Substances)~\cite{EU_Pesticides}, the OSPAR List of Chemicals for Priority Action~\cite{OSPAR_PriorityChemicals}, and a Flemish administrative substance list developed in the context of water quality modelling~\cite{Vlaanderen_SommatieStoffen}. This Flemish \texttt{summation substances} list reflects the heterogeneity of \emph{substance names} in the Flemish environmental permits and - monitoring. For each entry, available identifiers (CAS number, EC number, and substance name) were collected.

Where possible, substances were linked to structure-based identifiers (InChIKey) via PubChem~\cite{Kim2015}. The resulting dataset establishes a unified mapping between regulatory entries, their identifiers, and their originating source lists.

\paragraph{Chemical Identity and Linkability}
\label{sec:linkability}

Each substance entry was classified into one of several types reflecting its
degree of chemical specificity: structure-defined molecule, substance group,
mixture/UVCB, analytical parameter, or administrative entry.

Based on this classification, a linkability taxonomy was defined, ranging from
fully structure-resolvable entries (InChIKey) to non-identifiable entries.
This taxonomy expresses the extent to which regulatory substances can be
unambiguously linked across datasets.

\paragraph{Identifier Consistency Analysis}

For entries with both CAS numbers and InChIKeys, bidirectional mappings were
analysed to assess identifier consistency. This analysis evaluates whether CAS
numbers uniquely map to structures and vice versa.

\paragraph{Structure Coverage and Cross-List Overlap}

For each regulatory source, the proportion of structure-defined substances was
calculated. Overlap between sources was analysed using structure-based matching
(InChIKey) and set similarity measures.

\paragraph{Analysis of Non-Structure Substances}

Substances without structure identifiers were analysed using text-based
methods, including embedding-based clustering~\cite{Rousseeuw1987,Lloyd1982,McInnes2018} 
and similarity matching to chemical taxonomies (ChemOnt), to assess whether 
semantic techniques can recover meaningful structure or classification~\cite{Manning2008,DjoumbouFeunang2016}.
Each substance name and ChemOnt class label was encoded using the \texttt{all-MiniLM-L6-v2} sentence-transformer model~\cite{Reimers2019,Wang2020MiniLM,devlin-etal-2019-bert}, yielding a 384-dimensional dense vector representation.

\paragraph{Knowledge Graph Construction}

All data were integrated into an RDF knowledge graph using SKOS~\cite{W3C2009SKOS} for taxonomy
representation. Structure-defined substances were linked to ChemOnt classes and
to ChEBI entities via InChIKey, enabling integration of biological hazard
annotations~\cite{W3C2017OA}.

\paragraph{Workload Estimation}

For non-structure substances, interoperability requires manual assessment of
pairwise relationships (equivalence, subset, overlap, disjoint). The required
effort scales quadratically with the number of such entries as $\binom{N}{2} = N(N-1)/2$ where $N$ is the number of non-structure entries under consideration.

\paragraph{Priority Scoring}
\label{sec:priority_score}

To demonstrate the analytical capabilities of the integrated knowledge graph, a composite priority score was developed by integrating two independent dimensions: (i) regulatory presence across multiple European substance lists, including the CLP Regulation~\cite{CLP2008}, the EU Chemicals Strategy for Sustainability~\cite{EUCSS2020}, and the Persistent Organic Pollutants (POPs) Regulation~\cite{POPs2019}; and (ii) biological role annotations obtained from ChEBI~\cite{Degtyarenko2008}. The approach is conceptually grounded in the ToxPi (Toxicological Priority Index) framework~\cite{Reif2010}, which integrates heterogeneous evidence streams into a single prioritisation
score through transparent, domain-specific weights. 

This score is not intended to constitute a formal risk assessment. Rather, it serves as an illustrative aggregation of heterogeneous regulatory and biological information within a unified semantic framework. The objective is to demonstrate how interoperable, structure-linked data enable cross-domain queries and comparative analyses that are difficult to achieve using fragmented regulatory datasets. Moreover, such an approach can support the prioritisation of policy enforcement activities, particularly in contexts where regulatory agencies operate with limited personnel resources.

Overall, this methodology highlights the value of integrating regulatory metadata with ontology-based annotations to facilitate the identification of substances and chemical classes that are prominent across multiple dimensions of interest.